\documentclass[11pt,a4paper]{article}
% Projeto PTT6 — compilar com: tectonic main.tex  (ou latexmk -pdf main.tex)
\usepackage[T1]{fontenc}
\usepackage[utf8]{inputenc}
\usepackage[brazil]{babel}
\usepackage{lmodern}
\usepackage[a4paper,top=3cm,left=3cm,right=2cm,bottom=2cm]{geometry}
\usepackage{amsmath,amssymb}
\usepackage{float}
\usepackage{booktabs,tabularx,longtable,array,multirow}
\usepackage{enumitem}
\usepackage{xcolor}
\usepackage[hidelinks]{hyperref}
\usepackage[numbers,sort&compress]{natbib}
\usepackage{setspace}
\onehalfspacing
\setlist{noitemsep,topsep=2pt}
\newcolumntype{L}{>{\raggedright\arraybackslash}X}
\newcommand{\pend}[1]{\textcolor{orange!80!black}{[#1]}}

\title{\textbf{Benchmarking de Inferência do YOLO Quantizado em MCU e SBC}\\[2pt]
\large Desempenho e eficiência energética do YOLOX-Nano int8 para detecção de EPIs em \emph{Edge AI}\\[2pt]
\normalsize Projeto de Pesquisa --- Trilha PTT6 (protocolo experimental v1.5)}
\author{Filipe Matias \and Guilherme Lucas\\[2pt]
\small Orientador: Prof. Herbert Oliveira Rocha\\
\small Mestrado Profissional em Computação Aplicada (PPCA) --- Universidade Federal de Roraima\\
\small Disciplina: Arquitetura e Organização de Computadores}
\date{Outubro de 2026}

\begin{document}
\maketitle

\begin{abstract}
\noindent Este projeto compara, de forma reprodutível, a execução de um mesmo modelo de detecção de
objetos --- YOLOX-Nano quantizado em int8 --- em dois extremos do espectro \emph{Edge AI}: um
microcontrolador ESP32 (Xtensa LX6) e um computador de placa única Raspberry Pi Zero 2 W. A tarefa é a
detecção de Equipamentos de Proteção Individual (EPIs) no conjunto Ultralytics Construction-PPE,
começando pela classe colete e ampliando para as 11 classes. Seguindo a prática dos benchmarks de referência da área, a acurácia (mAP) é tratada como condição de
validade e reportada lado a lado com latência ponta a ponta (pré, inferência e pós-processamento),
throughput, energia por inferência --- estimada por software a partir de um modelo de potência por modo e
dos tempos medidos por fase ---, inferências por joule e EDP.
A hipótese H1 afirma que o SBC tem throughput maior, mas energia por inferência igual ou superior. O
desenho prevê piloto para definir repetições e warm-up, coleta em blocos independentes, estatística
robusta (Mann-Whitney U, $\delta$ de Cliff e IC por \emph{bootstrap}) e análise de sensibilidade do
modelo de potência.

\medskip\noindent\textbf{Palavras-chave:} Edge AI; TinyML; quantização int8; YOLOX; benchmarking; eficiência energética.
\end{abstract}

\section{Introdução}

A inferência de redes neurais na borda (\emph{Edge AI}) desloca o processamento da nuvem para
dispositivos próximos ao sensor, reduzindo latência, tráfego e exposição de dados. No outro lado,
impõe restrições severas de memória, capacidade de computação e energia. Entre as plataformas de borda,
dois extremos se destacam: \textbf{microcontroladores (MCU)}, de consumo muito baixo e memória da ordem
de centenas de kilobytes, e \textbf{computadores de placa única (SBC)}, que executam um sistema
operacional completo e oferecem mais desempenho a um custo energético maior.

O monitoramento do uso de EPIs em canteiros de obra é uma aplicação em que essa escolha é concreta: a
câmera precisa operar continuamente, frequentemente com alimentação restrita. Falta, porém, uma
comparação direta e reprodutível do \emph{mesmo} detector quantizado nas duas classes de plataforma.

\subsection{Pergunta de pesquisa e hipóteses}
\begin{description}
\item[Pergunta:] como se comparam latência, throughput e energia por inferência entre uma MCU (ESP32 LX6)
e um SBC (Raspberry Pi Zero 2 W) ao executar o mesmo YOLOX-Nano int8 na detecção de EPIs?
\item[H1:] o SBC apresenta throughput maior, mas energia por inferência igual ou superior, dado o seu
consumo estático mais alto.
\item[H0:] não há diferença nas métricas entre as plataformas.
\item[H2:] no SBC, a quantização int8 reduz latência e energia em relação ao fp32 ao custo de uma perda de
mAP, que é quantificada e reportada como trade-off.
\end{description}

\subsection{Objetivos}
\textbf{Geral:} comparar desempenho e eficiência energética da inferência do YOLOX-Nano int8 em uma MCU e
em um SBC.
\textbf{Específicos:}
\begin{enumerate}
\item preparar um artefato único \texttt{.tflite} int8 e validar sua equivalência funcional nas duas plataformas;
\item medir latência por fase, throughput e acurácia (mAP) sob um protocolo declarado;
\item estimar a energia por inferência por um modelo de potência parametrizado e quantificar sua sensibilidade;
\item testar H1 com estatística robusta e publicar os dados brutos (CSV) e o código.
\end{enumerate}

\subsection{Vínculo com a disciplina}
A comparação mobiliza desempenho e a Lei de Amdahl (tempos entre arquiteturas e o ganho de paralelismo
nos quatro núcleos do SBC), a hierarquia de memória (SRAM/PSRAM do ESP32 versus DRAM do SBC) e E/S.

\section{Fundamentação e trabalhos relacionados}
Benchmarks de inferência embarcada como o MLPerf Tiny~\cite{banbury2021mlperftiny} padronizam tarefas,
métricas (latência e energia) e regras de medição para microcontroladores, medindo apenas a
inferência. Noronha et al.~\cite{noronha2025} propõem integrar energia e latência separando três fases
(pré-inferência, inferência e pós-inferência) e adotam o EDP como métrica de eficiência --- abordagem
que este projeto segue na janela de medição. Fixa-se o runtime (TFLite int8) e varia-se apenas a
plataforma, para isolar o efeito do hardware.
Para a metodologia de repetições, seguimos a recomendação de definir o número de execuções a partir da
variabilidade observada, e não de um valor fixo~\cite{maricq2018}. O YOLOX~\cite{ge2021yolox} é um
detector \emph{anchor-free}; sua variante Nano tem $\approx$0,91\,M parâmetros e 1,08\,GFLOPs a
$416\times416$.


\paragraph{Métrica principal.} Na literatura de benchmarking de inferência, energia e latência não são
resultados válidos isoladamente: o MLPerf Inference associa a cada tarefa um \emph{quality target}
(tipicamente 99\% da acurácia FP32 de referência) que o sistema precisa atingir para que o desempenho conte
\cite{reddi2020mlperf}, e o MLPerf Tiny mede acurácia, latência e energia exigindo uma acurácia mínima por
modelo \cite{banbury2021mlperftiny}. Trabalhos de detecção voltados a energia reportam acurácia e energia
como eixos de um trade-off (fronteira de Pareto), não como um número composto
\cite{tu2023femtodet,deputter2025poq}, e o detector YOLO quantizado para MCU mais próximo do nosso caso
reporta mAP ao lado da energia por inferência \cite{moosmann2023tinyissimoyolo}; surveys de TinyML e de
quantização tratam a avaliação como multiobjetivo \cite{ray2022tinyml,gholami2021survey}. Adotamos,
portanto, acurácia como condição de validade e os demais eixos lado a lado.

\paragraph{Quantização.} A quantização int8 com aritmética inteira \cite{jacob2018quantization} bem
executada mantém a queda de mAP de detectores abaixo de cerca de 1 p.p. \cite{wu2020integer}; a avaliação
compara mAP int8 e FP32 no mesmo conjunto de teste \cite{nagel2021whitepaper}. Conjuntos de calibração da
ordem de uma centena de amostras costumam bastar para estimar faixas de ativação
\cite{krishnamoorthi2018quantizing}, mas ativações não-ReLU (sigmoide/Swish) sem suporte dedicado podem
degradar fortemente a acurácia quantizada \cite{nagel2021whitepaper}.

\paragraph{Orçamento de treino.} As 300 épocas do YOLOX referem-se a treino do zero no COCO
\cite{ge2021yolox}; a partir de pesos pré-treinados, o fine-tuning requer bem menos iterações
\cite{he2019rethinking}, e estudos de EPI com YOLO escolhem o número de épocas empiricamente, em geral com
parada antecipada \cite{otgonbold2022shel5k,ali2022labsafety}. Seguimos a parada pela validação
\cite{prechelt1998earlystopping} e as recomendações de relato e reprodutibilidade
\cite{dodge2019showyourwork,pineau2021reproducibility,henderson2018deeprl}.

\section{Materiais}

\subsection{Plataformas}
\begin{table}[H]\centering\small
\caption{Plataformas comparadas (fator independente).}\label{tab:plat}
\begin{tabularx}{\linewidth}{lLL}\toprule
 & \textbf{MCU} & \textbf{SBC}\\\midrule
Dispositivo & ESP32 (Xtensa LX6, dual-core) & Raspberry Pi Zero 2 W (Cortex-A53, 4 núcleos)\\
Clock & 240\,MHz, sem DFS/light-sleep & 1\,GHz, governor \texttt{performance}\\
Memória & 520\,KB SRAM + PSRAM externa \pend{confirmar módulo} & 512\,MB LPDDR2\\
Runtime & TensorFlow Lite Micro (int8) & TFLite/LiteRT (int8), 4 threads\\
Alimentação & USB 5\,V & fonte oficial 5\,V/2,5\,A\\\bottomrule
\end{tabularx}
\end{table}

A proposta original citava o ESP32-S3; adotou-se o ESP32 (LX6) disponível (emenda v1.1), sem as
extensões vetoriais do S3. Os resultados não se extrapolam para o S3.

\subsection{Modelo}
\emph{(v1.5)} O artefato avaliado até aqui é o YOLOX-Nano padrão (ativação SiLU), não a variante \emph{ti-lite}
prevista; a escolha final depende do diagnóstico de quantização (Seção~\ref{sec:quant}).

YOLOX-Nano na variante \emph{ti-lite} (amigável à quantização), treinado uma única vez por configuração
de classes no conjunto de treino (fora do escopo medido) e quantizado pós-treino em int8
\emph{full-integer} com 750 imagens de calibração estratificadas. O mesmo arquivo \texttt{.tflite}, com
hash SHA-256 registrado, é executado nas duas plataformas.

\subsection{Conjunto de dados}
Ultralytics Construction-PPE~\cite{ultralytics_cppe} (AGPL-3.0), com 11 classes: \emph{helmet},
\emph{gloves}, \emph{vest}, \emph{boots}, \emph{goggles}, \emph{none}, \emph{Person},
\emph{no\_helmet}, \emph{no\_goggle}, \emph{no\_gloves} e \emph{no\_boots}. O bloco A usa só a classe
\emph{vest}; o bloco B, as 11 classes. A Tabela~\ref{tab:ds} resume as contagens observadas no arquivo
baixado.

\begin{table}[H]\centering\small
\caption{Construction-PPE: imagens e instâncias da classe \emph{vest} por split.}\label{tab:ds}
\begin{tabular}{lrrr}\toprule
Split & Imagens & Instâncias \emph{vest} & Imagens com \emph{vest}\\\midrule
Treino & 1142 & 1283 & 837\\
Validação & 143 & 171 & 109\\
Teste & 141 & 178 & 111\\\bottomrule
\end{tabular}
\end{table}

\section{Metodologia experimental}\label{sec:metodo}

\subsection{Estratégia em fases}
O trabalho avança em fases com critério de saída: \textbf{F1} validação do modelo com imagens estáticas
(sem câmera); \textbf{F3} coleta das medições; \textbf{F4} análise e entrega. A integração da câmera ao vivo
(\textbf{F2}) passou a ser um \emph{extra} opcional (emenda v1.4): não bloqueia a medição nem entra no teste
de H1. Nenhuma medição final ocorre antes de a F1 estar funcional. Estado atual:
o modelo já executa no Raspberry Pi; a integração no ESP32 e a instrumentação final estão em andamento.

\subsection{Variáveis}
\begin{table}[H]\centering\small
\caption{Variáveis do experimento.}\label{tab:var}
\begin{tabularx}{\linewidth}{lL}\toprule
\textbf{Tipo} & \textbf{Variáveis}\\\midrule
Independente & Plataforma $\in\{$MCU, SBC$\}$\\
Dependentes & $t_{pre}$, $t_{inf}$, $t_{post}$, latência $t_{total}$ (ms); throughput (inf/s); energia por
inferência $E$ (mJ); inferências/J; EDP (mJ$\cdot$ms); mAP@50 e mAP@50-95\\
Controladas & artefato (hash), resolução comum $r$, quantização, limiares (conf.\ $\geq$0,25; IoU$_{NMS}$
0,45; até 100 detecções), mesmas 50 entradas, clock/governor, threads, rádio, alimentação, ambiente\\
Covariáveis & temperatura da CPU, frequência e \emph{throttling} (SBC), temperatura ambiente, nº de
detecções, imagem, bloco e ordem de execução\\\bottomrule
\end{tabularx}
\end{table}

\subsection{Métricas e critério de qualidade}\label{sec:metricas}
Para cada combinação \{plataforma $\times$ precisão\} reportam-se, lado a lado: mAP@50 e mAP@50:95 no
conjunto de teste (IC 95\% por \emph{bootstrap} de imagens \cite{efron1993bootstrap}), AP@50 do colete,
latência por fase, throughput, energia estimada, inferências/J e EDP. A retenção de acurácia
$\rho = \mathrm{mAP}_{\mathrm{int8}}/\mathrm{mAP}_{\mathrm{fp32}}$ é reportada como variável contínua; o alvo de
referência é $\rho \geq 0{,}99$ \cite{reddi2020mlperf}. Artefatos abaixo do alvo têm seus resultados de
energia rotulados como ``artefato degradado''. Como o mesmo artefato int8 roda nas duas plataformas, a
acurácia é idêntica entre elas (verificada no E1) e H1 compara plataformas sob acurácia controlada; H2
compara precisões no SBC como trade-off.

\subsection{Protocolo de treinamento}\label{sec:treino}
Fine-tuning a partir dos pesos COCO do YOLOX-Nano, com o número de épocas tratado como hiperparâmetro:
teto declarado, cronograma do YOLOX (warm-up; Mosaic/MixUp desligados nas últimas épocas) \cite{ge2021yolox},
mAP de validação a cada época, parada antecipada por paciência \cite{prechelt1998earlystopping} e checkpoint
final pelo melhor mAP de validação. Ao menos uma seed com curva completa (uma segunda se houver tempo,
reportando a faixa) \cite{henderson2018deeprl}; todos os hiperparâmetros são reportados
\cite{dodge2019showyourwork}. O valor fixo de 80 épocas usado na primeira rodada não tem respaldo e é
substituído por este protocolo.

\subsection{Avaliação de acurácia (E0) e quantização}\label{sec:quant}
O E0 avalia fp32 e int8 no conjunto de teste (141 imagens, 1.251 anotações) com o mesmo pré e
pós-processamento dos dispositivos, entrada em BGR (ordem de canais do treino), confiança 0,001 e NMS 0,65
(limiares de implantação 0,25/0,45 só para demonstração) e AP por classe. A calibração usa 750 imagens do
treino em BGR; pesos são quantizados por canal e ativações por tensor. A Tabela~\ref{tab:e0} mostra o E0 da
primeira rodada.

\begin{table}[H]\centering\small
\caption{E0 da primeira rodada (artefatos SiLU, calibração com 100 imagens em RGB).}\label{tab:e0}
\begin{tabular}{llrrr}\toprule
Artefato & Canais & Split & mAP@50 & mAP@50:95\\\midrule
fp32 & BGR & validação & 60,3\% & 29,1\%\\
fp32 & BGR & teste & 55,0\% & 27,2\%\\
fp32 & RGB & teste & 48,6\% & 23,3\%\\
int8 & BGR & teste & 16,7\% & 6,4\%\\
int8 & RGB (usado no RPi) & teste & 4,7\% & 1,5\%\\\bottomrule
\end{tabular}
\end{table}

O int8 retém $\rho \approx 0{,}30$ do mAP@50 do fp32 (BGR), muito abaixo do alvo e da faixa de PTQ bem
executada \cite{wu2020integer}; o colapso já está presente antes da concatenação da saída (17,1\% de
mAP@50 lendo os tensores internos), e a troca de canais custa 6,4 p.p. mesmo em fp32. O diagnóstico mínimo
compara (Q1) recalibração com 750 imagens BGR no mesmo grafo e (Q2) a variante com ReLU
\cite{nagel2021whitepaper}; segue para a medição final o artefato de maior $\rho$. As latências e energias já
medidas no SBC caracterizam o grafo SiLU, mas só serão apresentadas como desempenho de detector com o
artefato corrigido.

\subsection{Pipeline e janela de medição}
A latência cobre a janela completa \emph{pré + inferência + pós}, e as três fases são registradas
separadamente. Na fase estática, as entradas são tensores brutos $r\times r\times3$ uint8 já em
\emph{letterbox}, idênticos byte a byte nas duas plataformas; o pré-processamento medido é a cópia e a
quantização para o tensor de entrada. O pós-processamento é a decodificação da saída do YOLOX e o NMS
por classe, com a mesma lógica nas duas pontas. A resolução $r$ é comum: a maior de
$\{416, 320, 256, 224, 160\}$ que caiba na memória do ESP32 e permita ao menos 30 inferências por bloco
em até 60\,min.

\subsection{Ingestão direta de imagens}
Validação e medição usam as mesmas imagens do conjunto de teste, pré-processadas uma vez no PC e
\emph{ingeridas diretamente no modelo}. No ESP32 (core Arduino ou ESP-IDF com TensorFlow Lite Micro), o PC
envia cada imagem pela serial USB-UART em um quadro com cabeçalho e CRC32; o firmware a guarda em um buffer
em RAM, confirma o recebimento e só então abre a janela cronometrada; os tempos e as detecções voltam pela
serial depois da janela. No Raspberry Pi, as imagens são lidas do disco para a RAM antes do laço. Assim, a
janela começa com a imagem já em memória nas duas plataformas, e o custo de transferência fica fora da
medição. A serial foi preferida à flash porque 50 imagens a $160\times160$ ocupam $\approx$3,8\,MB, mais
que a partição de dados típica do ESP32.

\subsection{Instrumentação}
Os tempos vêm de relógios monotônicos de alta resolução em software (\texttt{time.perf\_counter\_ns()} no
SBC; \texttt{esp\_timer\_get\_time()} no MCU). Não há E/S dentro da janela: as medidas ficam em memória e
são gravadas entre blocos. Telemetria (temperatura, frequência, \emph{throttling}) é lida fora da
janela. Cada execução gera um CSV com uma linha por inferência e um arquivo de metadados (versões,
configurações, hashes, ambiente).

\subsection{Experimentos}
\begin{table}[H]\centering\small
\caption{Experimentos (repetidos para os blocos A --- 1 classe --- e B --- 11 classes).}\label{tab:exp}
\begin{tabularx}{\linewidth}{llL}\toprule
\textbf{Id} & \textbf{Fase} & \textbf{Descrição e saída}\\\midrule
E0 & F1 & Acurácia fp32 $\times$ int8 no teste (141 imagens), BGR, IC \emph{bootstrap}: mAP@50, mAP@50:95, AP do colete\\
E1 & F1 & Equivalência funcional: mesmas 50 entradas em PC, SBC e MCU; mesmas classes e nº de detecções, IoU $\geq$0,95, $|\Delta$score$|\leq$ 1 passo de quantização\\
E2 & F1/F3 & Piloto: viabilidade de $r$ no MCU; 100 inferências por plataforma; warm-up $W$ e $N$\\
E3 & F3 & Benchmark principal estático: 3 blocos por plataforma, $N=\max(30,N_{piloto})$\\
E3b & F3 & SBC com 1, 2 e 4 threads: fração paralelizável por Amdahl (exploratório)\\
E3c & F3 & SBC a $416\times416$ quando $r<416$ (exploratório)\\
E4 & F3 & Throughput sustentado em 60\,s de laço fechado e deriva térmica\\
E5 & extra & Câmera ao vivo (opcional, fora de H1): FPS sustentado e latência com captura\\
E6 & F4 & Energia derivada, sensibilidade e \emph{break-even}\\\bottomrule
\end{tabularx}
\end{table}

\paragraph{Warm-up e número de repetições.} Em janelas de 10 inferências, descartam-se as iniciais até a
variação relativa da mediana entre janelas consecutivas ficar abaixo de 2\% em duas janelas seguidas
(piso de 5). Com o coeficiente de variação (CV) de $t_{total}$ medido no piloto,
\begin{equation}
N=\left\lceil\left(\frac{1{,}96\,\mathrm{CV}}{\varepsilon}\right)^{2}\right\rceil,\qquad \varepsilon=1\%,\qquad N_{final}=\max(30,N),
\end{equation}
podendo $N$ diferir entre plataformas.

\paragraph{Estrutura de um bloco (E3).} Religar o dispositivo e aguardar 2\,min ocioso (registrar temperatura
e \emph{throttling}); executar $W$ inferências de warm-up (marcadas e excluídas); executar $\lceil N_{final}/3\rceil$
inferências cronometradas, percorrendo as 50 entradas em ordem fixa; gravar CSV e metadados; aguardar o
resfriamento (temperatura de volta a $\pm$2\,°C da inicial no SBC, ou 5\,min no MCU). Os blocos são
intercalados entre plataformas para diluir efeitos de dia e ambiente.

\subsection{Modelo de energia}
Como a medição física (INA219) foi descartada por aumentar demais a complexidade e o prazo, a energia é
estimada por software:
\begin{equation}
E_i = P_{pre}\,t_{pre,i}+P_{inf}\,t_{inf,i}+P_{post}\,t_{post,i},
\end{equation}
com a potência total do dispositivo em cada modo (Tabela~\ref{tab:pot}).

\begin{table}[H]\centering\small
\caption{Parâmetros do modelo de potência (nominal e faixa de sensibilidade).}\label{tab:pot}
\begin{tabularx}{\linewidth}{llrrL}\toprule
\textbf{Plataforma} & \textbf{Modo} & \textbf{Nominal} & \textbf{Faixa} & \textbf{Fonte}\\\midrule
SBC (placa, 5\,V) & ocioso & 0,60\,W & 0,50--1,00\,W & \cite{cnx2021zero2w,geerling2021zero2w}\\
SBC & ativo, 1 núcleo & 1,00\,W & 0,75--1,30\,W & \cite{geerling2021zero2w}\\
SBC & ativo, 4 núcleos & 2,30\,W & 1,85--2,90\,W & \cite{geerling2021zero2w,toms2021zero2w}\\
MCU (chip, 3,3\,V) & ativo, 240\,MHz & 0,224\,W & 0,099--0,224\,W & \cite{espressif_esp32_ds}\\
MCU & \emph{overhead} da placa & 0\,W & 0--0,25\,W & parâmetro declarado\\\bottomrule
\end{tabularx}
\end{table}

Com potência constante por modo, $E$ é uma transformação linear dos tempos; a incerteza do modelo é
tratada por sensibilidade, não pela estatística. A condição de \emph{break-even} é
\begin{equation}
E_{SBC}\geq E_{MCU}\iff \frac{P_{SBC}}{P_{MCU}}\geq\frac{t_{MCU}}{t_{SBC}},
\end{equation}
e reporta-se se a razão de potências necessária cai dentro da faixa plausível. Como os valores do SBC são
de placa e os do MCU de chip, H1 é avaliada também no pior caso para o MCU. Para um cenário de aplicação a
$f$ quadros/s, reporta-se ainda $E_{ciclo}(f)=P_{ativo}\,t+P_{ocioso}(1/f-t)$, válido para $t<1/f$.

\subsection{Análise estatística}
Por plataforma, reportam-se mediana e IQR, média $\pm$ desvio-padrão e CV das inferências válidas de E3.
A comparação usa o teste de Mann-Whitney U (bilateral para latência e unilateral nas direções de H1), o
$\delta$ de Cliff e intervalos de confiança de 95\% por \emph{bootstrap} BCa (10\,000 reamostras) para a
diferença e a razão das medianas. H1 é avaliada como teste de interseção--união, com
$\text{throughput}_{SBC}>\text{throughput}_{MCU}$ e $E_{SBC}\geq E_{MCU}$, cada parte a $\alpha=0{,}05$;
a parte energética é avaliada no valor nominal e nos extremos da Tabela~\ref{tab:pot}. Outliers não são
removidos; a remoção por Tukey (1,5$\times$IQR) entra apenas como sensibilidade. Verificam-se o efeito de
bloco (Kruskal-Wallis) e a autocorrelação (ACF de \emph{lag} 1). Uma análise secundária usa apenas
$t_{inf}$, isolando o núcleo de inferência das diferenças de implementação do pré e do pós-processamento.

\section{Ameaças à validade}
\begin{table}[H]\centering\small
\caption{Principais ameaças e mitigação.}\label{tab:ameacas}
\begin{tabularx}{\linewidth}{LL}\toprule
\textbf{Ameaça} & \textbf{Mitigação}\\\midrule
Energia estimada, não medida & mesmo método nas duas plataformas; premissas registradas; sensibilidade e \emph{break-even}\\
Escopos distintos (placa $\times$ chip) & \emph{overhead} de placa do MCU como parâmetro; H1 no pior caso para o MCU\\
Throttling térmico no SBC & governor fixo, dissipador, pausas; bloco com \emph{throttling} é repetido\\
Pré/pós em linguagens diferentes & mesma lógica e limiares; análise secundária só com $t_{inf}$\\
ESP32 LX6 em vez de S3 & desvio declarado; sem extrapolação para o S3\\
Resolução reduzida pela MCU & mesma $r$ nas duas plataformas; SBC a 416 apenas exploratório\\
Artefato int8 com baixa retenção de mAP & E0 obrigatório antes de E1--E4; energia rotulada ``degradado'' abaixo do alvo\\
Ordem de canais (BGR/RGB) inconsistente & BGR fixo em treino, calibração, avaliação e dispositivos\\
Ingestão por serial (MCU) $\times$ disco (SBC) & transferência fora da janela; janela inicia com a imagem em RAM\\
Câmera ao vivo (extra) & fora do teste de H1; se feita, mesma resolução, formato e nº de buffers\\
Autocorrelação e deriva & blocos independentes, ordem registrada, ACF e efeito de bloco\\\bottomrule
\end{tabularx}
\end{table}

\section{Cronograma e estado atual}
\begin{table}[H]\centering\small
\caption{Cronograma (15 semanas) e estado em outubro de 2026.}\label{tab:crono}
\begin{tabularx}{\linewidth}{llL}\toprule
\textbf{Semanas} & \textbf{Fase} & \textbf{Estado}\\\midrule
1--2 & Revisão e protocolo & concluída (protocolo v1.5; hardware e dataset definidos)\\
3--5 & Implementação & SBC medido (grafo SiLU); E0 revelou int8 degradado: retreino e requantização (v1.5)\\
6--9 & Experimentos & E2 (piloto) assim que $r$ e os artefatos forem fixados\\
10--11 & Análise & ---\\
12--13 & Relatório & ---\\
14--15 & Vídeo e entrega & ---\\\bottomrule
\end{tabularx}
\end{table}

\section{Entregáveis}
Repositório versionado com licença MIT; README de reprodutibilidade; dados brutos em CSV e metadados;
vídeo de demonstração (3--5\,min) nas duas plataformas, mostrando a ingestão de imagens e as detecções; relatório técnico conclusivo. Após a disciplina:
dataset com DOI (Zenodo) e submissão a evento regional.

\bibliographystyle{unsrtnat}
\bibliography{referencias}
\end{document}
